\documentclass{ximera}[font=12pt]  % I don't think font sizing works here

\title{Fractions: Simplifying}   % Use xmlatex all -s to compile when SERVE refuses to work.
\author{Kelly Stady}
\license{CC: 0}         % replace with an appropriate license, or set it in xmPreamble

\begin{document}


\begin{abstract}

\end{abstract}

\maketitle

\label{xim:FracsSimplify}


\section*{Simplifying Fractions}

A fraction is said to be in \textbf{simplest form} or \textbf{lowest terms} when the
numerator and denominator have no common factors other than 1.  For example, the fraction
$\textstyle \frac{2}{5}$ is in simplest form because 2 and 5 have no common factors.  However, the fraction
$\textstyle \frac{4}{6}$ is not in simplest form because both 4 and 6 have a common factor of 2.  The
process of writing a fraction in simplest form is called \textbf{simplifying} the fraction.  We will consider two methods 
to simplify fractions.  The first method uses the GCF and the second method uses prime factorizations.

\begin{procedure}[\textbf{Simplifying Fractions Using the GCF}]

    Divide the numerator and denominator by the GCF.

\end{procedure}

\begin{example}

Simplify the fraction $\textstyle \frac{4}{12}$ using the GCF = $4.$

\begin{explanation}

Divide the numerator and denominator by $4$ and write the result. 

\begin{align*}
\frac{4}{12} &= \frac{4 \div 4}{12 \div 4} \\[1em]
             &= \frac{\answer{1}}{\answer{3}}
\end{align*}

Done!

We can also simplify fractions by dividing the numerator and denominator by \emph{any} common factor, however using the GCF involves just one step.
If we didn't notice that the GCF was 4, but knew that both numbers were evenly divisible by 2, we would need two steps to simplify 
the fraction $\textstyle \frac{4}{12}.$  First, we divide the numerator and denominator by 2, then we divide by 2 again.  

\begin{align*}
\frac{4}{12} &= \frac{4 \div 2}{12 \div 2} \\[1em]
             &= \frac{\answer{2}}{\answer{6}}             \\[1em]
             &= \frac{2 \div 2}{6 \div 2} \\[1em]
             &= \frac{\answer{1}}{\answer{3}}
\end{align*}


\end{explanation}
\end{example}


\begin{problem}
  Simplify the fraction $\frac{24}{36}$ using the GCF = $12.$

  \[
  \frac{24}{36} = \frac{\answer{2}}{\answer{3}}
  \]

\end{problem}



\begin{problem}
  Simplify the fraction below.

  \[ \frac{18}{45} \]

  First, find the Greatest Common Factor (GCF) of 18 and 45.

  GCF = $\answer{9}$

  Now, divide both the numerator and the denominator by the GCF to simplify:

  \[ \frac{18}{45} = \frac{\answer{2}}{\answer{5}} \]

  \begin{hint}
    List the factors for each number to find the largest one they share (GCF):
    \begin{itemize}
        \item Factors of 18: 1, 2, 3, 6, 9, 18
        \item Factors of 45: 1, 3, 5, 9, 15, 45
    \end{itemize}
  \end{hint}
\end{problem}

\begin{procedure}[\textbf{Simplifying Fractions Using Prime Factorizations}]

Write the prime factorization of the numbers in the numerator and denominator, then cancel the common prime factors.

\end{procedure}

\begin{example}
  
Simplify the fraction $\textstyle \frac{20}{25}$ using prime factorizations.

\begin{explanation}

Write the prime factorization of each number then cancel the common factors and write the result. 

\begin{align*}
\frac{20}{25} &= \frac{2 \cdot 2 \cdot 5}{5 \cdot 5} \\[1em]
              &= \frac{2 \cdot 2 \cdot \cancel{5}}{5 \cdot \cancel{5}} \\[1em]
              &= \frac{4}{5}
\end{align*}


\end{explanation}
\end{example}


INSERT VIDEO of simplifying a fraction using one of these facts.

\begin{problem}
Simplify the fraction using the prime factorizations shown below.

\begin{eqnarray*}
    42 &=& 2 \cdot 3 \cdot 7 \\
    70 &=& 2 \cdot 5 \cdot 7
\end{eqnarray*}

$\displaystyle \frac{42}{70} = \frac{\answer{3}}{\answer{5}} $

  \begin{hint}
    Write the prime factorizations in the fraction, then cancel the common prime numbers.
    \[ \frac{42}{70} = \frac{2 \cdot 3 \cdot 7}{2 \cdot 5 \cdot 7} \]
  \end{hint}

\end{problem}


\begin{problem}
Simplify the fraction using whichever method you prefer.

\[ \frac{33}{60} = \frac{\answer{11}}{\answer{20}}\]

\begin{hint}
   If you want to use prime factorizations, write the prime factorizations in the fraction, then cancel the common prime numbers.

    \[ \frac{33}{60} = \frac{3 \cdot 11}{2 \cdot 2 \cdot 3 \cdot 5} \]

  \end{hint}

  \begin{hint}
    If you want to use the Greatest Common Factor (GCF), the largest number that divides both 33 and 60 is 3.
  \end{hint}

\end{problem}



\begin{problem}
  Simplify the fraction using whichever method you prefer.


  \[ \frac{60}{126} = \frac{\answer{10}}{\answer{21}} \]

  \begin{hint}
   If you want to use prime factorizations, write the prime factorizations in the fraction, then cancel the common prime numbers.

    \[ \frac{60}{126} = \frac{2 \cdot 2 \cdot 3 \cdot 5}{2 \cdot 3 \cdot 3 \cdot 7} \]

  \end{hint}

  \begin{hint}
    If you want to use the Greatest Common Factor (GCF), the largest number that divides both 60 and 126 is 6.
  \end{hint}

\end{problem}


\begin{problem}
  A recycling center collected $54$ tons of plastic last month.  However, due to contamination (like food waste or non-recyclable items), 
  only $21$ tons could be recycled.  What \textbf{fraction} of the collected plastic was recycled?  \textbf{Write your answer in simplest form}.

  \[ \text{Fraction Recycled} = \frac{\answer{7}}{\answer{18}} \]

  \begin{hint}
    Set up the fraction as the amount recycled over the total amount collected, then simplify.
    \[ \frac{21}{54} \]
  \end{hint}

\end{problem}


\begin{problem}
An animal rescue organization saved $120$ dogs from an overcrowded shelter.  Of these dogs, $45$ were seniors who required specialized 
medical care and longer stays before they were ready for adoption.  What \textbf{fraction} of the rescued dogs were seniors?  
\textbf{Write your answer in simplest form}.

\[ \text{Fraction of Senior Dogs} = \frac{\answer{3}}{\answer{8}} \]

\begin{hint}
Write the fraction as the number of senior dogs over the total number of dogs rescued, then simplify.

\[ \frac{45}{120} \]
\end{hint}
\end{problem}


\end{document}

